\subsection{\texorpdfstring{\(F_{4}\) 型系统}{F\_4 型系统}}

\begin{enumerate}
\def\labelenumi{(\Roman{enumi})}
\tightlist
\item
  考虑第 4 号的群 \(\mathbf{L}_2\)，它位于 \(\mathbf{R}^4\) 中。令 \(\mathbf{R}\) 为满足 \(\alpha \in \mathbf{L}_2\) 且 \((\alpha |\alpha) = 1\) 或 \((\alpha |\alpha) = 2\) 的向量集合；它包含向量
\end{enumerate}

\[
\pm \varepsilon_ {i}, \quad \pm \varepsilon_ {i} \pm \varepsilon_ {j} (i <   j), \quad \frac {1}{2} (\pm \varepsilon_ {1} \pm \varepsilon_ {2} \pm \varepsilon_ {3} \pm \varepsilon_ {4}).
\]

反之，若 \(\alpha\in R\)，则 \(\alpha\) 的坐标只能取 \(0,\pm\frac{1}{2},\pm1\) 这些值（因为 \((\frac{3}{2})^{2}>2\)）；这些坐标要么全为整数，从而得到向量 \(\pm\varepsilon_{i},\pm\varepsilon_{i}\pm\varepsilon_{j}\)，要么全都等于 \(\pm\frac{1}{2}\)，从而得到向量 \(\frac{1}{2}(\pm\varepsilon_{1}\pm\varepsilon_{2}\pm\varepsilon_{3}\pm\varepsilon_{4})\)。

我们证明，对于 \(\alpha, \beta \in \mathbb{R}\)，有 \(2(\alpha|\beta) / (\alpha|\alpha) \in \mathbf{Z}\)。若 \(\alpha = \pm \varepsilon_i\) 或 \(\alpha = \frac{1}{2} (\pm \varepsilon_1 \pm \varepsilon_2 \pm \varepsilon_3 \pm \varepsilon_4)\)，则 \((\alpha|\alpha) = 1\)，且由第 4 号知 \((\alpha|\beta) \in \frac{1}{2}\mathbf{Z}\)，因为 \(\alpha, \beta \in \mathbf{L}_2\)。若 \(\alpha = \pm \varepsilon_i \pm \varepsilon_j\)，则 \((\alpha|\alpha) = 2\)，且由第 4 号知 \((\alpha|\beta) \in \mathbf{Z}\)，因为 \(\alpha \in \mathbf{L}_1\) 且 \(\beta \in \mathbf{L}_2\)。因此，R 是 \(\mathbf{R}^4\) 中的约化根系（第 4 号）。根的数目为 \(n = 8 + \binom{4}{2} 4 + 2^4 = 48\)。

\begin{enumerate}
\def\labelenumi{(\Roman{enumi})}
\setcounter{enumi}{1}
\tightlist
\item
  在 \(\mathbf{R}^4\) 中赋予由基 \((\varepsilon_1, \varepsilon_2, \varepsilon_3, \varepsilon_4)\) 定义的字典序（§ 1，第 7 号）。特别地，\(\varepsilon_1 > \varepsilon_2 > \varepsilon_3 > \varepsilon_4\)。正根为
\end{enumerate}

\[
\varepsilon_ {i}, \quad \varepsilon_ {i} \pm \varepsilon_ {j} (i <   j), \quad \frac {1}{2} (\varepsilon_ {1} \pm \varepsilon_ {2} \pm \varepsilon_ {3} \pm \varepsilon_ {4}).
\]

最小的根是 \(\alpha_{3} = \varepsilon_{4}\)。在属于 \(\mathbf{R}\varepsilon_{3} + \mathbf{R}\varepsilon_{4}\) 但不属于 \(\mathbf{R}\varepsilon_{4}\) 的根中，最小的是 \(\alpha_{2} = \varepsilon_{3} - \varepsilon_{4}\)。在属于 \(\mathbf{R}\varepsilon_{2} + \mathbf{R}\varepsilon_{3} + \mathbf{R}\varepsilon_{4}\) 但不属于 \(\mathbf{R}\varepsilon_{3} + \mathbf{R}\varepsilon_{4}\) 的正根中，最小的是 \(\alpha_{1} = \varepsilon_{2} - \varepsilon_{3}\)。在不属于 \(\mathbf{R}\varepsilon_{2} + \mathbf{R}\varepsilon_{3} + \mathbf{R}\varepsilon_{4}\) 的正根中，最小的是 \(\alpha_{4} = \frac{1}{2} (\varepsilon_{1} - \varepsilon_{2} - \varepsilon_{3} - \varepsilon_{4})\)。没有一个 \(\alpha_{i}\) 是两个正根之和。因此，\((\alpha_{1},\alpha_{2},\alpha_{3},\alpha_{4})\) 是 R 的一个基（§ 1，第 6 号，命题 19 的推论 1）。我们有 \(\| \alpha_1 \|^2 = \| \alpha_2 \|^2 = 2\)，\(\| \alpha_3 \|^2 = \| \alpha_4 \|^2 = 1\)，\((\alpha_1|\alpha_2) = (\alpha_2|\alpha_3) = -1\)，\((\alpha_3|\alpha_4) = -\frac{1}{2}\)，\((\alpha_1|\alpha_3) = (\alpha_1|\alpha_4) = (\alpha_2|\alpha_4) = 0\)。可见 R 的 Dynkin 图为 \(F_4\) 型，因而不可约。

\begin{enumerate}
\def\labelenumi{(\Roman{enumi})}
\setcounter{enumi}{2}
\item
  我们有 \(h = \frac{n}{l} = 12\)。
\item
  令 \(\tilde{\alpha} = \varepsilon_1 + \varepsilon_2 = 2\alpha_1 + 3\alpha_2 + 4\alpha_3 + 2\alpha_4\)。\(\tilde{\alpha}\) 相对于 \((\alpha_i)\) 的坐标之和为 \(11 = h - 1\)，所以 \(\tilde{\alpha}\) 是最高根。我们有 \((\tilde{\alpha}|\alpha_1) = 1\)，\((\tilde{\alpha}|\alpha_2) = (\tilde{\alpha}|\alpha_3) = (\tilde{\alpha}|\alpha_4) = 0\)。
\end{enumerate}

扩展 Dynkin 图为

\[
\alpha_ {1} \quad \alpha_ {2} \quad \alpha_ {3} \quad \alpha_ {4}
\]

\begin{enumerate}
\def\labelenumi{(\Alph{enumi})}
\setcounter{enumi}{21}
\tightlist
\item
  公式 \(\alpha^{\prime} = \frac{2\alpha}{(\alpha|\alpha)}\) 给出 \(R^{\prime}\) 的向量集合 \(\pm 2\varepsilon_{i}, \pm \varepsilon_{i} \pm \varepsilon_{j}, \pm \varepsilon_{1} \pm \varepsilon_{2} \pm \varepsilon_{3} \pm \varepsilon_{4}\)。\(R^{\prime}\) 的 Dynkin 图按照第 2 号所说明的步骤由 \(R\) 的 Dynkin 图得到；于是我们看到 \(R^{\prime}\) 是 \(F_{4}\) 型的。
\end{enumerate}

与 \(\beta = \varepsilon_1\) 不正交的根为 \(\pm \varepsilon_1, \pm \varepsilon_1 \pm \varepsilon_j\)（ \(j \geqslant 2\)）以及 \(\frac{1}{2} (\pm \varepsilon_1 \pm \varepsilon_2 \pm \varepsilon_3 \pm \varepsilon_4)\)；数值 \(n(\alpha | \beta) = 2(\alpha | \beta)\) 对其中前 14 个根等于 \(\pm 2\)，对后 16 个根等于 \(\pm 1\)；因此，对于 \(\Phi_R\)，\(\beta\) 的长度平方为 \(4(14.4 + 16.1)^{-1} = \frac{1}{18}\)；于是

\[
\Phi_ {R} (x, y) = \frac {(x | y)}{18}.
\]

现在对 \(x = y = \beta\) 应用 § 1 第 12 号的公式 (18)，得到

\[
2 + 12. \frac {1}{4} + 16. \frac {1 / 4}{1} = \gamma (\mathrm{R}). \frac {1}{18}
\]

所以

\[
\gamma (\mathrm{R}) = 2. 3 ^ {4}.
\]

\begin{enumerate}
\def\labelenumi{(\Roman{enumi})}
\setcounter{enumi}{5}
\tightlist
\item
  计算基本权得到
\end{enumerate}

\[
\omega_ {1} = \varepsilon_ {1} + \varepsilon_ {2} = 2 \alpha_ {1} + 3 \alpha_ {2} + 4 \alpha_ {3} + 2 \alpha_ {4} = \tilde {\alpha}
\]

\[
\omega_ {2} = 2 \varepsilon_ {1} + \varepsilon_ {2} + \varepsilon_ {3} = 3 \alpha_ {1} + 6 \alpha_ {2} + 8 \alpha_ {3} + 4 \alpha_ {4}
\]

\[
\omega_ {3} = \frac {1}{2} (3 \varepsilon_ {1} + \varepsilon_ {2} + \varepsilon_ {3} + \varepsilon_ {4}) = 2 \alpha_ {1} + 4 \alpha_ {2} + 6 \alpha_ {3} + 3 \alpha_ {4}
\]

\[
\omega_ {4} = \varepsilon_ {1} = \alpha_ {1} + 2 \alpha_ {2} + 3 \alpha_ {3} + 2 \alpha_ {4}.
\]

\begin{enumerate}
\def\labelenumi{(\Roman{enumi})}
\setcounter{enumi}{6}
\tightlist
\item
  正根之和为
\end{enumerate}

\[
2 \rho = 11 \varepsilon_ {1} + 5 \varepsilon_ {2} + 3 \varepsilon_ {3} + \varepsilon_ {4} = 16 \alpha_ {1} + 30 \alpha_ {2} + 42 \alpha_ {3} + 22 \alpha_ {4}.
\]

\begin{enumerate}
\def\labelenumi{(\Roman{enumi})}
\setcounter{enumi}{7}
\item
  我们有 \(\mathrm{Q}(\mathrm{R}) = \mathrm{L}_2\)（第 4 号），且由 (VI) 有 \(\mathrm{P}(\mathrm{R}) = \mathrm{Q}(\mathrm{R})\)。因此，连接指标为 1。
\item
  指数族有 4 项，由于 \(h = 12\)，它必须包含与 12 互素的整数 1、5、7、11（§ 1，第 11 号，命题 30）；因此，这些就是 \(W(R)\) 的全部指数。
\item
  (X) 和 (XI) Dynkin 图的唯一自同构是恒等自同构，所以 \(A(R) = W(R)\) 且 \(w_{0} = -1\)。令 \(R'\) 为 R 的最长元所成的集合，即 \(\pm\varepsilon_{i} \pm \varepsilon_{j}\)；\(R'\) 是第 8 号构造的 \(D_{4}\) 型根系。显然，\(A(R)\) 的每个元素都是 \(A(R')\) 的元素。反之，\(A(R')\) 的元素保持 \(L_{1}\) 不变（因为它由 \(R'\) 生成），因而也保持与之关联的 \(L_{2}\) 不变，进而保持 R 不变。所以 \(W(R) = A(R) = A(R')\)。由第 8 号，\(W(R)\) 是 \(S_{3}\) 与 \(W(R')\) 的半直积，而 \(W(R')\) 本身是 \(S_{4}\) 与 \((\mathbf{Z}/2\mathbf{Z})^{3}\) 的半直积。\(W(R)\) 的阶为 \(3!4!2^{3} = 2^{7}.3^{2}\)。
\end{enumerate}

